\documentclass[11pt,a4paper]{article}
\usepackage[margin=25mm]{geometry}
\usepackage{amsmath,amssymb,amsthm,hyperref}
\setlength{\emergencystretch}{2em}
\newtheorem{theorem}{Theorem}
\newtheorem{lemma}[theorem]{Lemma}
\newcommand{\F}{\mathbf F}
\newcommand{\gen}[1]{\langle#1\rangle}
\newcommand{\Om}{\Omega}
\title{Six negative reflections in an orthogonal space of dimension five}
\author{Human Compression project: F001 working component}
\date{}
\begin{document}\maketitle
We treat one of the two involution types in odd dimension five. This
argument does not exclude strings containing an involution whose
negative eigenspace has dimension two. That remaining case is required
for the full upper bound in Brooksbank's rank spectrum theorem.
This draft has not received distinct review.

\section{The classification input}
Let $V$ be a nondegenerate orthogonal space of dimension five over
$\F_q$, where $q$ is odd, and write $B$ for its symmetric bilinear
form. The reflection with root $v$, where $B(v,v)\ne0$, is
\[
 s_v(x)=x-2\frac{B(x,v)}{B(v,v)}v.
\]
We retain the following part of the finite reflection classification.
A finite irreducible orthogonal group in dimension five, generated by
five reflections, is, after extending scalars and conjugating, either
\begin{enumerate}
\item $\mathrm O_5(q_0)$ or one of its subgroups generated by one
reflection class, for a finite field of odd order $q_0$ of the same
characteristic; or
\item a faithful modular reduction of the standard reflection group
$A_5$, $B_5$, or $D_5$.
\end{enumerate}
The groups in the first case contain $\Om_5(q_0)$ and have a
nondegenerate invariant form over $\F_{q_0}$. Allowing any $q_0$
avoids assuming that the originally chosen coefficient field is
minimal. This is the dimension-five specialization of
Monson--Schulte~\cite[Theorem 3.1]{MS}, restated with its proof
in Monson's author-hosted notes~\cite[Theorem 4.1, pp.15--17]{M}.
Those notes explicitly exclude the deleted symmetric module and the
exceptional modular $E_6$ representation when the number of generating
reflections is only the dimension. The underlying classification is
due to Zalesskii--Serezhkin~\cite{ZS}. Its full proof remains an
external acquisition cost.

We also retain the classical order formula
\[
 \omega(t):=|\Om_5(t)|=\tfrac12t^4(t^2-1)(t^4-1),
 \qquad |\mathrm O_5(t)|=4\omega(t).                 \tag{1}
\]
We retain also the simplicity of $\Om_5(q)$ for odd $q\ge5$.
The homomorphism $\pi:\mathrm O(V)\longrightarrow\mathrm{SO}(V)$,
$g\mapsto\det(g)g$, has kernel $\{1,-1\}$.

\section{The end groups and their field of definition}
\begin{lemma}\label{path}
Suppose $s_0,\ldots,s_{r-1}$ are distinct reflections with roots
$v_0,\ldots,v_{r-1}$, commuting when their indices differ by more
than one and not commuting at consecutive indices. Then $r\le6$.
When $r=6$, either consecutive set of five roots is a basis, and
each corresponding end reflection group acts absolutely irreducibly.
\end{lemma}
\begin{proof}
Distinct orthogonal reflections commute precisely when their roots
are perpendicular. The Gram matrix is tridiagonal with nonzero
entries immediately above and below the diagonal. A kernel vector
is determined by its first coordinate: the first row determines the
second coordinate, the second row the third, and so on. A nonzero
kernel vector has both end coordinates nonzero, by applying this
argument from either end. Thus its rank is at least $r-1$ and at
most five.

For $r=6$ the rank is five. The roots span $V$ and their
one-dimensional relation space is the Gram kernel. No relation
omitting an end root is possible, so either end five are a basis.
Let $H$ be one end group and extend scalars to an algebraic closure.
If $0\ne W\le V$ is $H$-invariant, either every reflection fixes
$W$ pointwise, or some root lies in $W$, since
$(s_i-1)W\subseteq\gen{v_i}$. The first case places $W$ in the
perpendicular of a basis, which is zero. In the second, consecutive
nonzero pairings propagate membership along the path, placing all
five basis roots in $W$. Hence $W=V$.
\end{proof}

\begin{lemma}\label{descent}
Suppose $r=6$ above and $\gen{-s_0,\ldots,-s_5}=\Om_5(q)$.
If $H=\gen{s_0,\ldots,s_4}$ is of orthogonal type in the
retained classification, then $\pi(H)=\Om_5(q)$.
\end{lemma}
\begin{proof}
Conjugate over an algebraic closure so that $H$ and its invariant
nondegenerate form $B_0$ are defined over $\F_{q_0}$. Absolute
irreducibility shows that the original form $B$ is a scalar multiple
of $B_0$: the quotient of their matrix maps commutes with $H$.
For $i\le4$, the negative eigenspace of $s_i$ has a root defined
over $\F_{q_0}$. The first four root lines therefore span a
four-space defined over that field. Since $s_5$ commutes with
$s_0,\ldots,s_3$, its root line is exactly the perpendicular of
this four-space. It too is defined over $\F_{q_0}$, so $s_5$ is
defined there and preserves $B_0$. Consequently
\[
                 \gen{s_0,\ldots,s_5}\le\mathrm O_5(q_0)
                                                        \tag{2}
\]
after the same conjugation.

Since $H$ contains $\Om_5(q_0)$ and a reflection,
$|H|\ge2\omega(q_0)$. Also $\pi(H)\le\Om_5(q)$ and the
kernel of $\pi$ has order two, so $\omega(q_0)\le\omega(q)$.
The function $\omega(t)$ is strictly increasing for $t\ge3$,
whence $q_0\le q$. If $q_0<q$, these are powers of the same odd
prime, so $q\ge3q_0$. Formula (1) gives
$\omega(q)>4\omega(q_0)$, contradicting (2) and the order of
the signed image. Thus $q_0=q$. Finally
$|\pi(H)|\ge|H|/2\ge\omega(q)$, giving the assertion.
\end{proof}

\section{The finite alternatives}
\begin{lemma}\label{finite}
If an end group is a finite alternative, it has type $A_5$ or
$B_5$. Its four consecutive product orders are respectively
$(3,3,3,3)$ or $(4,3,3,3)$, up to reversal.
\end{lemma}
\begin{proof}
In the standard $A_5$ representation, an involution is a reflection
precisely when it is a transposition of the six coordinates. Its
faithful irreducible reduction has characteristic different from
two and three, and negative eigenspace dimensions are unchanged.
Five transpositions generate $S_6$ only when their edges form a
connected graph on six letters. With five edges it is a tree.
Consecutive transpositions do not commute, so each consecutive
product has order three.

In the signed permutation representation of $B_5$, the reflecting
root lines are the short lines $\gen{e_i}$ and the long lines
$\gen{e_i\pm e_j}$. This holds in every odd characteristic:
an involutory signed permutation has one negative eigenvalue for
each positive signed two-cycle and each negative fixed coordinate.
The long reflections are precisely the reflections in $D_5$.

For five reflections to generate $B_5$, their permutation images
must generate $S_5$. At least four are long. At least one is short,
since long reflections lie in $D_5$. There are therefore exactly
four long roots and one short root. The four long-root edges form
a tree on the five coordinates. If the short root $e_i$ were
internal in the string, its two neighboring long roots would both
involve $i$. Their edges are distinct and have only that coordinate
in common, so those roots would not be perpendicular, contradicting
their distance two in the string. The short root is at an end.
Its adjacent product has order four; adjacent long-root products
have order three.

In $D_5$ every noncommuting pair of reflections has product order
three. A connected string of five such reflections is a quotient
of the Coxeter group $A_5\cong S_6$, of order $720$, and cannot
generate $D_5$, of order $1920$.
\end{proof}

\begin{theorem}\label{bound}
Let $q\ge5$ be odd. A string C-group marking of $\Om_5(q)$
whose generators are all negatives of reflections has rank at most
five. More strongly, any connected six-term string of distinct negative
reflections generating $\Om_5(q)$ has at least one end group equal
to the whole group.
\end{theorem}
\begin{proof}
The diagram of a C-marking of this simple group is connected:
a disconnected string would give nontrivial commuting normal factors.
Lemma~\ref{path} therefore bounds its rank by six. Suppose six generators
exist and both signed end groups are proper. Lemma~\ref{descent}
excludes orthogonal type for either ordinary end group.
Lemma~\ref{finite} applies to both ends. Comparing their overlapping
four-bond sequences shows that the middle three bonds are all three;
each outer bond is either three or four.

Over an algebraic closure rescale the roots to have norm one.
Put $x_i=B(v_{i-1},v_i)^2$ for $1\le i\le5$.
On the span of two independent roots the reflection product has
determinant one and trace $4x_i-2$. Hence order three gives
$x_i=1/4$, including characteristic three, and order four gives
$x_i=1/2$. If $D_j$ is the leading $j$ by $j$ Gram determinant,
\[
 D_0=D_1=1,\qquad D_{j+1}=D_j-x_jD_{j-1}.          \tag{3}
\]
The four possibilities give
\[
\begin{array}{c|c}
(3,3,3,3,3)&D_6=7/64\\
(4,3,3,3,3),\ (3,3,3,3,4)&D_6=1/32\\
(4,3,3,3,4)&D_6=0.
\end{array}                                                   \tag{4}
\]
As $D_6=0$, one outer four is impossible. The all-three case has
characteristic seven and is a quotient of $A_6\cong S_7$. Its
signed image has order at most $5040<\omega(q)$.

For the remaining case compare the following six root lines in
the standard diagonal five-space:
\[
 e_1,\quad e_1-e_2,\quad e_2-e_3,\quad e_3-e_4,
 \quad e_4-e_5,\quad e_5.                           \tag{5}
\]
After normalizing norms over the algebraic closure, their squared
consecutive pairings agree with ours. Changing signs successively
makes the Gram matrices identical. The first five roots form a
basis with nondegenerate Gram determinant $1/16$. The isometry
mapping these five normalized roots to (5) maps the sixth as well,
since its pairings with a basis determine it. All six ordinary
reflections are therefore conjugate to signed permutation matrices.
Their group has order at most $|B_5|=3840<\omega(q)$, another
contradiction.

Thus one signed end group is full. In a C-marking each generator
is outside the subgroup on all the other marks, by the intersection
property. Neither end can be full, excluding rank six.
\end{proof}

\paragraph{Scope and attribution.}
The reflection classification and classical order formula are retained
results. This argument combines the field of definition of an end
group with the overlapping finite end types. The proper signed
$S_7$ counterexample to the source's arbitrary-subgroup lemma is
consistent with this theorem. No bound for mixed involution types,
complete F001 recovery, conjecture resolution, or priority is claimed.

\begin{thebibliography}{9}
\bibitem{MS} B.~Monson and E.~Schulte, \emph{Reflection groups and
polytopes over finite fields, I}, Advances in Applied Mathematics
\textbf{33} (2004), 290--317, doi:10.1016/j.aam.2003.11.002.
\bibitem{M} B.~Monson, \emph{Constructions of Symmetric Polytopes},
author-hosted course notes, supplied 54-page edition,
\url{https://bmonson.ext.unb.ca/fields/outline.pdf}, Section~4,
Theorem~4.1 and Table~3, pp.15--19.
\bibitem{ZS} A.~E.~Zalesskii and V.~N.~Serezhkin,
\emph{Finite linear groups generated by reflections},
Mathematics of the USSR-Izvestiya \textbf{17} (1981), 477--503.
\end{thebibliography}
\end{document}
